\chapter{振动与波}
\label{chap:oscillations-waves}

振动与波是高中物理中最能体现“微分方程思想”的主题之一。一个质点如果总被某种回复作用拉回平衡位置，就会产生振动；许多相邻质点之间若存在耦合，局部振动便会向外传播，形成波。用微积分的语言说，振动研究的是时间函数的微分方程，波研究的是时空函数的偏微分方程。

本章将以最典型的二阶线性方程为主线：先从简谐运动出发，再讨论单摆、阻尼振动和受迫振动，最后进入波动方程、干涉与驻波，并用经典题说明微积分视角如何把“记结论”变成“会推导”。

\section{简谐运动}

\subsection{从牛顿第二定律到微分方程}

设质点在一条直线上运动，平衡位置取为 $x=0$。若回复力满足胡克定律
\[
F=-kx,
\]
其中 $k>0$ 为劲度系数，则由牛顿第二定律 $m\ddot x=F$ 得
\[
 m\ddot x=-kx.
\]
整理为
\begin{equation}
\ddot x+\omega^2 x=0,
\label{eq:sho-ode}
\end{equation}
其中
\[
\omega\coloneqq \sqrt{\frac{k}{m}}
\]
称为\emph{角频率}。

\begin{definition}
满足微分方程
\[
\ddot x+\omega^2x=0
\]
的运动称为\emph{简谐运动}。其中振幅记为 $A$，周期记为 $T$，频率记为 $f$，相位记为 $\omega t+\varphi$，初相记为 $\varphi$。
\end{definition}

\begin{figure}[H]
\centering
\begin{tikzpicture}[scale=1]
  \draw[fill=gray!20] (-0.5,-0.8) rectangle (0,0.8);
  \draw[thick] (0,0) -- (0.5,0);
  \draw[thick] (0.5,0) -- (0.8,0.25) -- (1.1,-0.25) -- (1.4,0.25) -- (1.7,-0.25) -- (2.0,0.25) -- (2.3,-0.25) -- (2.6,0);
  \draw[fill=blue!15, thick] (2.6,-0.45) rectangle (4.1,0.45);
  \draw[->, thick] (3.35,0) -- (4.35,0) node[right] {$x$};
  \draw[dashed] (3.35,-0.7) -- (3.35,0.7);
  \node[below] at (3.35,-0.7) {平衡位置};
  \node at (3.35,0) {$m$};
  \node[above] at (1.55,0.35) {$k$};
\end{tikzpicture}
\caption{弹簧振子的理想模型}
\end{figure}

\subsection{通解与运动学量}

方程 \eqref{eq:sho-ode} 的特征方程为
\[
r^2+\omega^2=0,
\]
其根为 $r=\pm i\omega$。因此有如下结论。

\begin{theorem}
简谐运动方程
\[
\ddot x+\omega^2x=0
\]
的实数通解为
\begin{equation}
 x(t)=A\cos(\omega t+\varphi),
\label{eq:sho-general}
\end{equation}
其中 $A\ge 0$，$\varphi$ 为常数。等价地，也可写成
\[
 x(t)=C_1\cos \omega t+C_2\sin \omega t.
\]
其速度与加速度分别为
\[
 v(t)=\dot x=-A\omega\sin(\omega t+\varphi),
\qquad
 a(t)=\ddot x=-\omega^2x.
\]
故加速度始终指向平衡位置，且与位移成正比、反向。
\end{theorem}

由 $\omega=2\pi/T=2\pi f$ 可得
\[
T=\frac{2\pi}{\omega},
\qquad
f=\frac{\omega}{2\pi}.
\]
这说明只要知道系统的“恢复快慢”，就能确定振动节奏。

\begin{figure}[htbp]
\centering
\begin{tikzpicture}
\begin{axis}[
  width=0.75\textwidth, height=0.5\textwidth,
  xlabel={时间 $t$}, ylabel={位移 $x$},
  axis lines=middle,
  xmin=0, xmax=6.4, ymin=-1.2, ymax=1.2,
  grid=major, grid style={gray!30},
  thick,
]
\addplot[blue, domain=0:6.283, samples=160] {cos(deg(2*x + 0.6))};
\end{axis}
\end{tikzpicture}
\caption{简谐运动的位移随时间按余弦规律变化}
\end{figure}

\begin{remark}
把 $x=A\cos(\omega t+\varphi)$ 投影到圆周运动上理解，是学习简谐运动的经典方法：设一点在半径为 $A$ 的圆上做匀速圆周运动，其在某一直径上的投影便做简谐运动。这个几何图像能帮助我们理解相位差、超前与滞后。
\end{remark}

\subsection{能量观点}

简谐振动虽然速度和位移都在不断变化，但总机械能保持不变。

\begin{law}
对理想弹簧振子，动能与势能分别为
\[
K=\frac12 m\dot x^2,
\qquad
U=\frac12 kx^2=\frac12 m\omega^2x^2.
\]
总机械能
\begin{equation}
E=K+U=\frac12 m\dot x^2+\frac12 m\omega^2x^2
\label{eq:sho-energy}
\end{equation}
恒为常数，且
\[
E=\frac12 kA^2=\frac12 m\omega^2A^2.
\]
\end{law}

证明很直接：对式 \eqref{eq:sho-energy} 对时间求导，得
\[
\dv{E}{t}=m\dot x\ddot x+m\omega^2x\dot x
= m\dot x(\ddot x+\omega^2x)=0.
\]
因此 $E$ 守恒。

\begin{example}
已知某振子满足 $x(0)=3\text{ cm}$，$v(0)=-4\omega\text{ cm/s}$。求其振幅和一个可取的解析式。

\textbf{解：} 设
\[
x=C_1\cos\omega t+C_2\sin\omega t.
\]
由初始条件得
\[
C_1=3,
\qquad
v(0)=\omega C_2=-4\omega,
\]
故 $C_2=-4$，于是
\[
x=3\cos\omega t-4\sin\omega t.
\]
振幅为
\[
A=\sqrt{3^2+(-4)^2}=5\text{ cm}.
\]
若写成振幅—初相形式，则可取
\[
x=5\cos(\omega t+\varphi),
\qquad
\cos\varphi=\frac35,
\quad
\sin\varphi=\frac45.
\]
因此可取 $\varphi=\arctan(4/3)$。
\end{example}

\begin{example}
一质量为 $m$ 的物体做简谐运动，振幅为 $A$，角频率为 $\omega$。求其最大速度、最大加速度，以及当 $x=A/2$ 时的速度大小。

\textbf{解：} 由
\[
v=-A\omega\sin(\omega t+\varphi)
\]
可知最大速度为
\[
v_{\max}=A\omega.
\]
由 $a=-\omega^2x$ 得最大加速度为
\[
a_{\max}=\omega^2A.
\]
当 $x=A/2$ 时，由机械能守恒
\[
\frac12 m v^2+\frac12 m\omega^2\left(\frac{A}{2}\right)^2
=\frac12 m\omega^2A^2,
\]
解得
\[
v=\pm \frac{\sqrt3}{2}A\omega.
\]
若问速度大小，则答案为 $\frac{\sqrt3}{2}A\omega$。
\end{example}

\section{单摆}

\subsection{单摆方程与小角近似}

长度为 $\ell$、摆球质量为 $m$ 的理想单摆，在竖直平面内摆动。取偏角 $\theta$ 为广义坐标，则切向方向上有
\[
m\ell\ddot\theta=-mg\sin\theta,
\]
即
\begin{equation}
\ddot\theta+\frac{g}{\ell}\sin\theta=0.
\label{eq:pendulum-exact}
\end{equation}
这就是单摆的精确方程。

\begin{figure}[H]
\centering
\begin{tikzpicture}[scale=1.15]
  \coordinate (O) at (0,2.2);
  \coordinate (A) at (0,0);
  \coordinate (B) at (1.3,0.32);
  \draw[thick] (-0.8,2.2) -- (0.8,2.2);
  \fill (O) circle (2pt);
  \draw[dashed] (O) -- (A);
  \draw[thick] (O) -- (B);
  \fill[blue!15] (B) circle (0.16);
  \pic [draw, ->, "$\theta$", angle eccentricity=1.35, angle radius=0.48cm] {angle = A--O--B};
  \node[right] at (0.7,1.1) {$\ell$};
  \draw[->] (B) -- ++(0,-1.0) node[below] {$mg$};
  \node[below] at (A) {平衡位置};
\end{tikzpicture}
\caption{单摆的几何模型}
\end{figure}

当偏角很小时，$\sin\theta\approx\theta$，于是 \eqref{eq:pendulum-exact} 近似为
\begin{equation}
\ddot\theta+\frac{g}{\ell}\theta=0.
\label{eq:pendulum-small-angle}
\end{equation}
它正是简谐运动方程，因此
\[
\theta(t)=\Theta_0\cos\left(\sqrt{\frac{g}{\ell}}\,t+\varphi\right),
\qquad
T\approx 2\pi\sqrt{\frac{\ell}{g}}.
\]

\begin{definition}
若用 $\sin\theta\approx\theta$ 将单摆方程线性化，则称为\emph{小角近似}。通常当 $\theta\lesssim 10^\circ$ 时，该近似对周期的误差已相当小；当振幅增大时，周期会比小角公式稍长。
\end{definition}

\subsection{精确周期与椭圆积分}

单摆不是严格的简谐振动。要得到精确周期，可从机械能出发。若摆角最大值为 $\theta_0$，则在任意时刻有
\[
\frac12 m\ell^2\dot\theta^2+mg\ell(1-\cos\theta)=mg\ell(1-\cos\theta_0).
\]
化简得
\[
\dot\theta^2=\frac{2g}{\ell}(\cos\theta-\cos\theta_0).
\]
从 $\theta=0$ 摆到 $\theta=\theta_0$ 所需时间是四分之一周期，于是
\[
\frac{T}{4}=\sqrt{\frac{\ell}{2g}}\int_0^{\theta_0}\frac{\dd\theta}{\sqrt{\cos\theta-\cos\theta_0}}.
\]
作代换
\[
\sin\frac{\theta}{2}=k\sin\phi,
\qquad
k=\sin\frac{\theta_0}{2},
\]
可化为
\begin{equation}
T=4\sqrt{\frac{\ell}{g}}\,K(k),
\qquad
K(k)=\int_0^{\pi/2}\frac{\dd\phi}{\sqrt{1-k^2\sin^2\phi}}.
\label{eq:elliptic-period}
\end{equation}
其中 $K(k)$ 称为\emph{第一类完全椭圆积分}。

\begin{remark}
式 \eqref{eq:elliptic-period} 已经给出了单摆周期的精确表达式。它不像简谐运动那样能用初等函数写出，但“无法写成初等函数”并不意味着不能研究：通过积分表达式，我们仍能得到近似、数值计算和误差估计。
\end{remark}

当 $\theta_0$ 不大时，椭圆积分可展开为
\[
T=2\pi\sqrt{\frac{\ell}{g}}\left(1+\frac{\theta_0^2}{16}+\frac{11\theta_0^4}{3072}+\cdots\right),
\]
这清楚说明：摆幅越大，周期越长。

\begin{example}
长度 $\ell=1.0\text{ m}$ 的单摆在某地做小振幅摆动，取 $g=9.8\text{ m/s}^2$。求其小角近似周期；若最大摆角为 $0.40\text{ rad}$，用二阶修正估计真实周期。

\textbf{解：} 小角近似下
\[
T_0=2\pi\sqrt{\frac{\ell}{g}}\approx 2.01\text{ s}.
\]
若考虑振幅修正，则
\[
T\approx T_0\left(1+\frac{\theta_0^2}{16}\right)
=2.01\left(1+\frac{0.40^2}{16}\right)
\approx 2.03\text{ s}.
\]
可见摆幅为 $0.40\text{ rad}$ 时，周期虽然只增加约百分之一，但已经偏离“与振幅无关”的理想结论。
\end{example}

\section{阻尼振动}

现实系统中总会有摩擦、空气阻力或内耗。若阻力近似与速度成正比，即 $f=-b\dot x$，则弹簧振子的运动方程变为
\[
m\ddot x+b\dot x+kx=0.
\]
除以 $m$ 后写成标准形式
\begin{equation}
\ddot x+2\gamma \dot x+\omega_0^2x=0,
\label{eq:damped}
\end{equation}
其中
\[
\gamma\coloneqq \frac{b}{2m},
\qquad
\omega_0\coloneqq \sqrt{\frac{k}{m}}.
\]

\begin{definition}
方程
\[
\ddot x+2\gamma\dot x+\omega_0^2x=0
\]
描述的振动称为\emph{阻尼振动}。参数 $\gamma$ 描述阻尼强弱，$\omega_0$ 是无阻尼时的固有角频率。
\end{definition}

令 $x=e^{rt}$，得到特征方程
\[
r^2+2\gamma r+\omega_0^2=0.
\]
其根为
\[
r=-\gamma\pm\sqrt{\gamma^2-\omega_0^2}.
\]

\begin{theorem}
阻尼振动按 $\gamma$ 与 $\omega_0$ 的大小关系分为三类：
\begin{enumerate}
  \item \textbf{欠阻尼}（$\gamma<\omega_0$）：此时根为共轭复根，解为
  \[
  x(t)=A e^{-\gamma t}\cos(\omega t+\varphi),
  \qquad
  \omega=\sqrt{\omega_0^2-\gamma^2}.
  \]
  物体仍振动，但振幅按指数规律衰减。
  \item \textbf{临界阻尼}（$\gamma=\omega_0$）：此时根为二重根 $r=-\gamma$，解为
  \[
  x(t)=(C_1+C_2 t)e^{-\gamma t}.
  \]
  物体不振荡，并且在“不发生往复”的前提下最快回到平衡位置附近。
  \item \textbf{过阻尼}（$\gamma>\omega_0$）：此时根为两个不同实根，解为
  \[
  x(t)=C_1 e^{r_1 t}+C_2 e^{r_2 t},
  \qquad
  r_{1,2}=-\gamma\pm\sqrt{\gamma^2-\omega_0^2}.
  \]
  物体同样不振荡，但回到平衡位置较慢。
\end{enumerate}
\end{theorem}

对欠阻尼情况，包络线为 $\pm A e^{-\gamma t}$。由于机械能与振幅平方成正比，故能量近似满足
\[
E(t)\propto e^{-2\gamma t}.
\]
这说明轻微阻尼下，振幅衰减得慢，能量却按更快的指数规律流失。

\begin{remark}
高中题里常说“振幅逐渐减小”，微积分语言给出更精确的说法：如果阻力与速度成正比，则振幅不是线性减小，而是指数衰减。也就是说，前几次看上去衰减明显，越到后面越“拖尾”。
\end{remark}

\begin{figure}[htbp]
\centering
\begin{tikzpicture}
\begin{axis}[
  width=0.75\textwidth, height=0.5\textwidth,
  xlabel={时间 $t$}, ylabel={位移 $x$},
  axis lines=middle,
  xmin=0, xmax=12, ymin=-1.1, ymax=1.1,
  grid=major, grid style={gray!30},
  thick,
]
\addplot[blue, domain=0:12, samples=220] {exp(-0.18*x)*cos(deg(4*x))};
\addplot[red, dashed, domain=0:12, samples=120] {exp(-0.18*x)};
\addplot[red, dashed, domain=0:12, samples=120] {-exp(-0.18*x)};
\end{axis}
\end{tikzpicture}
\caption{欠阻尼振动及其指数衰减包络线}
\end{figure}

\begin{example}
已知阻尼方程
\[
\ddot x+6\dot x+25x=0.
\]
判断阻尼类型，并写出其通解。

\textbf{解：} 与标准形式比较得
\[
2\gamma=6\Rightarrow \gamma=3,
\qquad
\omega_0=5.
\]
由于 $\gamma<\omega_0$，故为欠阻尼。阻尼后的角频率为
\[
\omega=\sqrt{\omega_0^2-\gamma^2}=\sqrt{25-9}=4.
\]
因此通解为
\[
x(t)=A e^{-3t}\cos(4t+\varphi).
\]
若改写为正弦余弦线性组合，则也可写成
\[
x(t)=e^{-3t}(C_1\cos4t+C_2\sin4t).
\]
\end{example}

\section{受迫振动与共振}

\subsection{受迫振动的稳态解}

若系统除恢复力与阻力外，还受到周期性驱动力
\[
F(t)=F_0\cos\Omega t,
\]
则方程变为
\begin{equation}
\ddot x+2\gamma\dot x+\omega_0^2x=\frac{F_0}{m}\cos\Omega t.
\label{eq:forced}
\end{equation}
它的总解由“瞬态解 $+$ 稳态解”组成。瞬态部分会因阻尼逐渐消失，长期保留下来的是与外力同频的稳态振动。

设稳态解形如
\[
x_p=X\cos(\Omega t-\delta),
\]
代入 \eqref{eq:forced} 可得振幅与相位差分别为
\begin{equation}
X(\Omega)=\frac{F_0/m}{\sqrt{(\omega_0^2-\Omega^2)^2+4\gamma^2\Omega^2}},
\label{eq:forced-amplitude}
\end{equation}
以及
\begin{equation}
\tan\delta=\frac{2\gamma\Omega}{\omega_0^2-\Omega^2}.
\label{eq:forced-phase}
\end{equation}
当 $\Omega$ 很小时，系统几乎跟着外力同步摆动，$\delta\approx 0$；当 $\Omega$ 很大时，位移明显滞后于驱动力，$\delta$ 接近 $\pi$。

\begin{definition}
在周期外力作用下，系统经过足够长时间后保留下来的与外力同频振动，称为\emph{稳态受迫振动}。若在某个驱动频率附近振幅显著增大，则称系统发生了\emph{共振}。
\end{definition}

\subsection{共振频率}

共振并不是“频率相同”这句话的机械重复，而是振幅函数 \eqref{eq:forced-amplitude} 取得最大值的问题。令分母最小，即研究
\[
D(\Omega)=(\omega_0^2-\Omega^2)^2+4\gamma^2\Omega^2.
\]
对 $\Omega$ 求导并令其为零，可得（除去 $\Omega=0$ 的平凡情形）
\[
\Omega_r=\sqrt{\omega_0^2-2\gamma^2}.
\]
因此只有当 $\gamma<\omega_0/\sqrt2$ 时，振幅曲线才在正频率处有明显峰值。

\begin{remark}
无阻尼极限 $\gamma\to 0$ 时，$\Omega_r\to\omega_0$，这就是教材常说“驱动力频率等于固有频率时共振”。一旦存在阻尼，共振峰会变矮、变宽，并且峰值位置略小于 $\omega_0$。
\end{remark}

\begin{figure}[htbp]
\centering
\begin{tikzpicture}
\begin{axis}[
  width=0.75\textwidth, height=0.5\textwidth,
  xlabel={驱动角频率 $\Omega/\omega_0$}, ylabel={稳态振幅 $X$},
  axis lines=middle,
  xmin=0, xmax=1.8, ymin=0, ymax=5.4,
  grid=major, grid style={gray!30},
  thick,
]
\addplot[blue, domain=0.05:1.8, samples=220] {1/sqrt((1 - x^2)^2 + 0.04*x^2)};
\addplot[gray, dashed] coordinates {(0.9899,0) (0.9899,5.2)};
\end{axis}
\end{tikzpicture}
\caption{阻尼受迫振动的共振曲线在固有频率附近出现峰值}
\end{figure}

\begin{example}
某受迫振动满足
\[
\ddot x+0.40\dot x+4x=2\cos\Omega t.
\]
求稳态振幅 $X(\Omega)$，并求共振频率。

\textbf{解：} 由标准形式可知
\[
2\gamma=0.40\Rightarrow \gamma=0.20,
\qquad
\omega_0=2,
\qquad
\frac{F_0}{m}=2.
\]
由 \eqref{eq:forced-amplitude}
\[
X(\Omega)=\frac{2}{\sqrt{(4-\Omega^2)^2+0.16\Omega^2}}.
\]
共振频率为
\[
\Omega_r=\sqrt{\omega_0^2-2\gamma^2}
=\sqrt{4-2\times 0.2^2}
=\sqrt{3.92}
\approx 1.98.
\]
它略小于无阻尼时的固有角频率 $2$。
\end{example}

\section{波动方程}

\subsection{从离散振子到连续介质}

波可以看成“振动的传播”。最典型的模型是张紧弦上的横波。设弦的线密度为 $\mu$，张力大小近似恒为 $T$，挠度记为 $y(x,t)$。取弦上一小段 $[x,x+\Delta x]$，在小振幅近似下，竖直方向合力约为
\[
T\pdv{y}{x}(x+\Delta x,t)-T\pdv{y}{x}(x,t)
\approx T\pdv[2]{y}{x}\Delta x.
\]
而这小段弦的质量为 $\mu\Delta x$，故由牛顿第二定律得
\[
\mu\Delta x\pdv[2]{y}{t}=T\pdv[2]{y}{x}\Delta x.
\]
约去 $\Delta x$，便得到一维波动方程
\begin{equation}
\pdv[2]{y}{t}=v^2\pdv[2]{y}{x},
\qquad
v\coloneqq \sqrt{\frac{T}{\mu}}.
\label{eq:wave}
\end{equation}

\begin{definition}
若函数 $y(x,t)$ 满足
\[
\pdv[2]{y}{t}=v^2\pdv[2]{y}{x},
\]
则称其描述了一维介质中的\emph{波动}，$v$ 称为\emph{波速}。
\end{definition}

\begin{theorem}
一维波动方程 \eqref{eq:wave} 的一般解可写为
\begin{equation}
 y(x,t)=F(x-vt)+G(x+vt),
\label{eq:dalembert}
\end{equation}
其中 $F,G$ 为任意二次可导函数。$F(x-vt)$ 表示沿 $+x$ 方向传播的行波，$G(x+vt)$ 表示沿 $-x$ 方向传播的行波。
\end{theorem}

验证并不困难。例如对 $y=F(x-vt)$，有
\[
\pdv{y}{t}=-vF'(x-vt),
\qquad
\pdv[2]{y}{t}=v^2F''(x-vt),
\]
而
\[
\pdv{y}{x}=F'(x-vt),
\qquad
\pdv[2]{y}{x}=F''(x-vt),
\]
所以确有 $\partial_t^2 y=v^2\partial_x^2 y$。

最常见的简谐行波写成
\begin{equation}
 y(x,t)=A\cos(kx-\omega t+\varphi),
\label{eq:harmonic-wave}
\end{equation}
其中
\[
\omega=vk,
\qquad
\lambda=\frac{2\pi}{k},
\qquad
f=\frac{\omega}{2\pi},
\qquad
v=f\lambda.
\]

\begin{figure}[H]
\centering
\begin{tikzpicture}[scale=0.95]
  \draw[->] (-0.2,0) -- (6.6,0) node[right] {$x$};
  \draw[->] (0,-1.6) -- (0,1.6) node[above] {$y$};
  \draw[thick, blue!70!black, domain=0:6.1, samples=120] plot (\x,{0.9*sin(deg(1.6*\x))});
  \draw[->, red!70!black, thick] (4.8,1.1) -- (5.8,1.1) node[right] {$v$};
  \node[above] at (2.7,1.0) {某时刻的行波形状};
\end{tikzpicture}
\caption{向右传播的简谐行波示意图}
\end{figure}

例如，函数
\[
y(x,t)=0.02\cos(4\pi x-200\pi t)
\]
就是式 \eqref{eq:harmonic-wave} 的具体实例，其中 $k=4\pi$、$\omega=200\pi$，故
\[
v=\frac{\omega}{k}=50,
\qquad
\lambda=\frac{2\pi}{k}=0.50,
\qquad
f=\frac{\omega}{2\pi}=100.
\]
它是 $F(x-vt)$ 型函数，所以满足一维波动方程，并沿 $+x$ 方向传播。

\section{波的叠加与干涉}

\subsection{叠加原理}

线性波动方程最重要的性质之一是线性叠加。

\begin{law}
若 $y_1(x,t)$ 与 $y_2(x,t)$ 都满足线性波动方程
\[
\pdv[2]{y}{t}=v^2\pdv[2]{y}{x},
\]
则任意线性组合
\[
y(x,t)=c_1y_1(x,t)+c_2y_2(x,t)
\]
仍然满足该方程。这称为\emph{叠加原理}。
\end{law}

叠加原理意味着不同波形可以“穿过彼此”而瞬时相加。宏观上看到的干涉图样，本质上正是相位关系造成的线性相加结果。

\subsection{驻波}

两列振幅相同、频率相同、波速相同、传播方向相反的简谐波
\[
y_1=A\cos(kx-\omega t),
\qquad
y_2=A\cos(kx+\omega t)
\]
叠加后得到
\begin{equation}
 y=2A\cos kx\cos\omega t.
\label{eq:standing-wave}
\end{equation}
这就是\emph{驻波}。它看起来“不向前传播”，因为空间部分 $\cos kx$ 与时间部分 $\cos\omega t$ 已经分离。

由 \eqref{eq:standing-wave} 可知：
\begin{itemize}
  \item 当 $\cos kx=0$ 时，振幅恒为零，称为\emph{波节}；
  \item 当 $|\cos kx|=1$ 时，振幅最大为 $2A$，称为\emph{波腹}。
\end{itemize}

若弦长为 $L$，两端固定，则边界条件要求
\[
L=n\frac{\lambda}{2}
\qquad (n=1,2,3,\dots),
\]
于是允许出现一系列本征频率
\[
f_n=n\frac{v}{2L}.
\]

\subsection{拍频}

若两列同方向传播的简谐波频率很接近，例如
\[
y_1=A\cos\omega_1 t,
\qquad
y_2=A\cos\omega_2 t,
\]
则叠加后
\begin{equation}
 y=2A\cos\frac{\omega_1-\omega_2}{2}t\cdot \cos\frac{\omega_1+\omega_2}{2}t.
\label{eq:beat}
\end{equation}
快速振动由平均角频率控制，缓慢变化的包络由差频控制。若记 $f_1=\omega_1/(2\pi)$、$f_2=\omega_2/(2\pi)$，则拍频为
\[
f_{\text{beat}}=|f_1-f_2|.
\]

\begin{remark}
拍频现象很适合用来调音：当两个音叉频率非常接近时，人耳会听到声音强弱周期性起伏。起伏越慢，说明两者频率越接近；当拍频消失时，两者几乎同频。
\end{remark}

\begin{example}
一根两端固定的弦长为 $L=0.60\text{ m}$，波速为 $120\text{ m/s}$。求其基频；再求频率为 $256\text{ Hz}$ 与 $260\text{ Hz}$ 的两列声波叠加时的拍频。

\textbf{解：} 两端固定弦的基频为
\[
f_1=\frac{v}{2L}=\frac{120}{1.2}=100\text{ Hz}.
\]
若考虑两声波拍频，则
\[
f_{\text{beat}}=|260-256|=4\text{ Hz}.
\]
这意味着每秒听到 $4$ 次强弱变化。
\end{example}

\begin{figure}[htbp]
\centering
\begin{tikzpicture}
\begin{axis}[
  width=0.75\textwidth, height=0.5\textwidth,
  xlabel={位置 $x$}, ylabel={位移 $y$},
  axis lines=middle,
  xmin=0, xmax=6.4, ymin=-2.2, ymax=2.2,
  grid=major, grid style={gray!30},
  thick,
  legend style={at={(0.98,0.98)}, anchor=north east, fill=white, draw=gray!30},
]
\addplot[blue, dashed, domain=0:6.283, samples=180] {cos(deg(x - 0.8))};
\addplot[red, dashed, domain=0:6.283, samples=180] {cos(deg(x + 0.8))};
\addplot[black, domain=0:6.283, samples=180] {cos(deg(x - 0.8)) + cos(deg(x + 0.8))};
\legend{$y_1$, $y_2$, $y_1+y_2$}
\end{axis}
\end{tikzpicture}
\caption{相向传播的两列简谐波叠加后形成驻波}
\end{figure}

下面选取三道常见振动、波动题型，说明微积分方法如何使题目的结构更清晰。

\section*{本章小结}

本章的核心线索可以概括为以下几点：
\begin{itemize}
  \item 简谐运动由二阶常系数线性方程 $\ddot x+\omega^2x=0$ 描述，其解为余弦函数，机械能守恒；
  \item 单摆的小角近似把非线性方程化为简谐方程，而精确周期需要借助椭圆积分；
  \item 阻尼振动对应方程 $\ddot x+2\gamma\dot x+\omega_0^2x=0$，可分为欠阻尼、临界阻尼和过阻尼；
  \item 受迫振动的稳态振幅取决于驱动频率，阻尼存在时共振峰出现于 $\Omega_r=\sqrt{\omega_0^2-2\gamma^2}$；
  \item 波动方程 $\partial_t^2 y=v^2\partial_x^2 y$ 的基本解是行波，叠加后可出现干涉、驻波和拍频。
\end{itemize}

从“振动”到“波”，物理图景在变，但数学思想始终统一：先建立微分方程，再结合初始条件、边界条件和对称性分析其解的结构。这正是微积分物理学的基本范式。
